\documentclass[11pt,a4paper]{book}
\usepackage[utf8]{inputenc}
\usepackage[T1]{fontenc}
% \usepackage{CJKutf8}
\usepackage{ctex}
\usepackage{amsmath,amssymb,amsthm,mathtools}
\usepackage{graphicx}
\usepackage{float}
\usepackage[dvipsnames,hyperref]{xcolor}
\usepackage{tikz}
\usetikzlibrary{angles,quotes,arrows.meta}
\usepackage{import}
\usepackage[colorlinks=true,linkcolor=MidnightBlue,citecolor=OliveGreen,urlcolor=RoyalBlue]{hyperref}
\usepackage[margin=1in]{geometry}
\usepackage{setspace}
\usepackage{lipsum}
\setstretch{1.05}

% Keep Chinese text support, but use English names for floats.
\ctexset{figurename=Figure, tablename=Table}

\newtheorem{definition}{Definition}[chapter]
\newtheorem{example}{Example}[chapter]
\newtheorem{note}{Note}[chapter]
\newtheorem{theorem}{Theorem}[chapter]

% Each chapter sets this once; module imports then stay local to that chapter.
\newcommand{\currentchapterpath}{}
\newcommand{\chapterinput}[1]{\subimport{\currentchapterpath}{#1.tex}}

\title{Mathematical Methods for Physics}
\author{Kai Luo}
\date{updated on \today}

\begin{document}

\frontmatter
\maketitle
\chapter*{Foreword}
\addcontentsline{toc}{chapter}{Foreword}

This book is written as a flexible lecture-note template for mathematical physics and related topics. The aim is to keep each chapter self-contained, while still allowing the full text to compile as a single coherent volume.

The structure is intentionally modular: each chapter sits in its own directory and can be developed independently. This makes it easier to draft, revise, and compile one chapter at a time without losing the overall book organization.

When a concept is first introduced it is set in italic type; an acronym is set in sans-serif type at its first appearance. An index of these terms, together with a list of references, appears at the end of the volume. Standard treatments of the same material include Arfken, Weber, and Harris~\cite{arfken}, Riley, Hobson, and Bence~\cite{riley}, and Boas~\cite{boas}.

The material in these pages is meant to be expanded, corrected, and adapted to your own teaching and research goals.

\tableofcontents

\mainmatter

% Compile only selected chapters at a time.
% Uncomment the next line and list the chapter files you want to build.
% \includeonly{chapters/prelim/prelim,chapters/vec/vec,chapters/fs/fs,chapters/ft/ft}

\chapter{Preliminaries}
\renewcommand{\currentchapterpath}{chapters/prelim/}

\begin{flushright}
\textit{“合抱之木，生于毫末；九层之台，起于累土。” ——《道德经》
\footnote{From the \textit{Dao De Jing}: ``A great tree grows from a tiny shoot; a nine-story terrace begins with a heap of earth.'' The quotation emphasizes that substantial achievements begin with small, steady steps.}}
\end{flushright}

This introduction chapter reviews some of the basic concepts and notions used
throughout the book. It aims to establish common vocabularies
and notations that are used in modern mathematical physics literature.
Though you, as a reader, may share little connection with me (the author),
it is hoped that after going through this chapter,
you and me are now connected in some way and we are on the same page.

A few notational conventions are used uniformly below.
The imaginary unit is $\mathrm{i}$ and the differential is $\mathrm{d}$,
both in upright roman type.
Vectors are written in bold, as in $\mathbf{A}$.
The Laplacian is $\nabla^2$; a finite increment is $\Delta x$.
A string or interval length is $\ell$.
The Fourier transform uses the symmetric convention
$F(\omega)=\mathcal{F}\{f\}(\omega)=(2\pi)^{-1/2}\int f(x)e^{-\mathrm{i}\omega x}\,\mathrm{d}x$.
Real logarithms are written $\ln$; the complex logarithm is $\log$.


 

\chapterinput{sets_and_logic}
\chapterinput{trigonometry}
\chapterinput{calculus}
\chapterinput{linear_algebra}

\include{chapters/vec/vec}
\include{chapters/complex/complex}
\include{chapters/pde/pde}
\chapter{Fourier Series}
\renewcommand{\currentchapterpath}{chapters/fs/}

\begin{flushright}
\textit{“八音克谐，无相夺伦。” ——《尚书·舜典》}
\footnote{From the ``Canon of Shun'' in the \textit{Book of Documents}: ``The eight classes of musical sounds achieve harmony, with none displacing the proper order of another.'' The line praises harmony formed from distinct components. Fourier analysis similarly represents a complicated signal as a superposition of elementary oscillations, each with its own frequency and amplitude; their ordered combination reconstructs the original behavior.}
\end{flushright}

Fourier series resolve periodic functions into elementary oscillations. They occur throughout signal processing, acoustics, optics, and the solution of differential equations on finite intervals. This chapter develops the orthogonal-function setting and the Fourier theorem for periodic functions~\cite{fourier,bracewell,lighthill}. The passage from a discrete trigonometric expansion to an integral transform is taken up in the next chapter.

\chapterinput{fourier_series}
\chapterinput{exercises}

\chapter{Fourier Transform and Laplace Transform}
\renewcommand{\currentchapterpath}{chapters/ft/}

\begin{flushright}
\textit{“往来不穷谓之通。” ——《周易·系辞上》}
\footnote{From the ``Appended Statements'' of the \textit{Book of Changes}: ``Coming and going without exhaustion is called penetration.'' An integral transform sends a function into another domain and the inverse returns it. Taken together, the two maps make a problem transparent: differentiation becomes multiplication, convolution becomes a product, and an initial-value problem becomes algebra.}
\end{flushright}

The Fourier transform extends a trigonometric expansion from periodic functions to functions on the real line. The Laplace transform is a closely related tool for causal functions and initial-value problems. This chapter passes from Fourier series to the Fourier integral, develops the delta function as the natural language of point sources and inversion, and then introduces the Laplace transform.

\chapterinput{fourier_transform}
\chapterinput{delta_function}
\chapterinput{laplace_transform}
\chapterinput{applications}
\chapterinput{exercises}

\include{chapters/cov/cov}
\chapter{Green's Function Method}
\renewcommand{\currentchapterpath}{chapters/gf/}

\begin{flushright}
\textit{“举一隅不以三隅反，则不复也。” ——《论语·述而》
\footnote{From the \textit{Analects}, ``Shu Er'': ``If I raise one corner and [the student] cannot return with the other three, I do not repeat myself.'' A Green's function is that one corner: the response to a unit point source. Once it is known, superposition returns the other three---the field of an arbitrary source under the same boundary and initial conditions.}}
\end{flushright}

Earlier chapters solved definite problems mainly by separation of variables. The \term{Green's function method} takes a complementary route. A Green's function---also called a \term{point-source influence function}---is the field produced by a unit point source subject to prescribed boundary and/or initial conditions. Once that response is known, linearity yields the field of an arbitrary source by superposition.

This chapter follows the classical development for Poisson's equation, the method of images, time-dependent wave and heat problems, the impulse theorem, and (optionally) generalized Green formulas for elliptic, parabolic, and hyperbolic operators~\cite{green1828,duffy,stakgold,jackson}.

\chapterinput{poisson_green}
\chapterinput{images}
\chapterinput{time_green}
\chapterinput{impulse}
\chapterinput{generalized}
\chapterinput{exercises}


\backmatter

\end{document}
